\chapter{L'omeostasi}

\textbf{Obiettivi del corso}: conoscere l'organizzazione funzionale del corpo umano; approfondire
il funzionamento dei diversi organi e i meccanismi generali di controllo funzionale; conoscere gli
adattamenti delle funzioni vitali dell'organismo in risposta all'attività sportiva.

% ---------------------------------------------------------------------------
\section{Definizione e concetto}

\textbf{Omeostasi}: mantenimento di condizioni stabili dell'organismo (equilibrio corporeo, stato
di benessere); insieme di processi e meccanismi fisiologici che correggono tempestivamente le
perturbazioni del sistema.
\begin{itemize}
  \item le condizioni non sono fisse $\rightarrow$ \textbf{equilibrio dinamico}: stabilità e
    bilanciamento, non immobilità;
  \item ``il problema vero della fisiologia è la conoscenza dei meccanismi di regolazione
    dell'organismo''.
\end{itemize}

\rubrica{Ambiente esterno e ambiente interno}
\begin{itemize}
  \item \textbf{ambiente esterno}: ambiente in cui vive l'organismo;
  \item \textbf{ambiente interno}: ambiente in cui vivono gli elementi che formano l'organismo
    (plasma e liquidi corporei).
\end{itemize}

\textbf{Regolazione delle interazioni tra più organi}: insieme di processi e meccanismi fisiologici
che correggono tempestivamente le perturbazioni del sistema.
\begin{itemize}
  \item \textbf{feedback}: risposta correttiva a una perturbazione già avvenuta;
  \item \textbf{feedforward}: concetto dinamico che tiene conto degli adattamenti anche predittivi
    $\rightarrow$ atteggiamenti comportamentali, regolazione anticipatoria.
\end{itemize}

% ---------------------------------------------------------------------------
\section{Ambiente interno e compartimenti liquidi}

\textbf{Scambi continui} tra i vari ambienti: attraverso l'ambiente extracellulare avvengono gli
scambi di energia e sostanze con il mondo esterno.
\begin{itemize}
  \item \textbf{organi/apparati interfaccia} (con l'esterno): cute, polmoni, reni, apparato
    gastroenterico.
\end{itemize}

\rubrica{Compartimenti dei liquidi corporei}
\begin{itemize}
  \item \textbf{liquido extracellulare (LEC)}: 20\% del peso corporeo;
    \begin{itemize}
      \item \textit{plasma sanguigno}: 5\% del peso corporeo;
      \item \textit{liquido interstiziale}: 15\% del peso corporeo;
    \end{itemize}
  \item \textbf{liquido intracellulare (LIC)}: 40\% del peso corporeo (citoplasma).
\end{itemize}

\textbf{Parametri da mantenere costanti}: pressione arteriosa, glicemia, pH, temperatura corporea,
composizione e distribuzione dei liquidi e degli elettroliti.

% ---------------------------------------------------------------------------
\section{Meccanismi di controllo}

Due modalità di intervento:
\begin{itemize}
  \item \textbf{a feedback}: la variazione di un parametro dal valore fisiologico, prodotta da una
    certa causa, induce una reazione che agisce sulla causa stessa; entra in azione \textit{solo
    dopo} che il cambiamento è avvenuto;
  \item \textbf{a feedforward}: interviene per controllare un processo \textit{prima} che avvenga
    la variazione (es.\ comando centrale nell'esercizio fisico);
    \begin{itemize}
      \item atteggiamenti comportamentali;
      \item regolazione anticipatoria.
    \end{itemize}
\end{itemize}

\rubrica{Controllo a feedback: componenti}
Il sistema di controllo a feedback è costituito da 3 elementi:
\begin{enumerate}
  \item \textbf{sensori} (recettori);
  \item \textbf{centro di integrazione};
  \item \textbf{effettori}.
\end{enumerate}

\rubrica{Tipi di feedback}
\begin{itemize}
  \item \textbf{feedback negativo}: la risposta si oppone alla variazione, riportando il parametro
    verso il valore desiderato;
  \item \textbf{feedback positivo}: (citato dal docente tra i tipi di controllo a feedback).
\end{itemize}

\rubrica{Logica del controllo omeostatico}
Per una \textbf{variabile controllata} (pH, pressione arteriosa, temperatura, ecc.):
\begin{itemize}
  \item il valore rientra nell'\textbf{intervallo desiderato} $\rightarrow$ nessun intervento;
  \item il valore \textbf{non} rientra nell'intervallo desiderato $\rightarrow$ attivazione del
    sensore $\rightarrow$ centro di integrazione $\rightarrow$ effettori;
  \item esito: \textbf{compensazione riuscita} $\rightarrow$ benessere; \textbf{compensazione
    fallita} $\rightarrow$ malattia.
\end{itemize}

% ---------------------------------------------------------------------------
\section{Esempio: la termoregolazione}

Mantenimento della temperatura corporea interna al valore fisiologico di $37\,^\circ$C $\pm$ 0,5.
\begin{itemize}
  \item \textbf{sensori}: termocettori centrali e periferici;
  \item \textbf{centro di controllo}: ipotalamo, confronta la temperatura rilevata con quella del
    \textit{set point} ($37\,^\circ$C $\pm$ 0,5);
  \item \textbf{effettori}: ghiandole sudoripare (attivazione), vasodilatazione;
  \item tipo di controllo: \textbf{feedback negativo}.
\end{itemize}

\rubrica{Risposte secondo la temperatura}
\begin{multicols}{2}
\textbf{Temperatura bassa} ($< 37\,^\circ$C $\pm$ 0,5) --- minimizzare la perdita e massimizzare la
produzione di calore:
\begin{itemize}
  \item vasocostrizione periferica;
  \item mancanza di sudorazione;
  \item risposte comportamentali per trattenere il calore;
  \item aumento della produzione di calore tramite il metabolismo;
  \item brivido;
  \item risposte comportamentali (p.\,es.\ muoversi).
\end{itemize}

\columnbreak

\textbf{Temperatura alta} ($> 37\,^\circ$C $\pm$ 0,5) --- massimizzare la perdita e minimizzare la
produzione di calore:
\begin{itemize}
  \item vasodilatazione periferica;
  \item sudorazione;
  \item risposte comportamentali per disperdere calore;
  \item diminuzione dell'apporto di cibo (riduce la produzione di calore metabolico);
  \item risposte comportamentali (p.\,es.\ riduzione dell'attività fisica).
\end{itemize}
\end{multicols}

% ---------------------------------------------------------------------------
\section{Omeostasi ed esercizio fisico}

\subsection{Adattamento e ipercompensazione}

\textbf{Adattamento}: condizioni fisiche e capacità funzionali migliori. Piccoli aumenti di
affaticamento e debolezza in seguito all'esercizio fisico sono seguiti da \textbf{ipercompensazione}
durante il recupero.
\begin{itemize}
  \item con allenamento: esercizio $\rightarrow$ debilitazione $\rightarrow$ ripristino
    $\rightarrow$ adattamento $\rightarrow$ (con ulteriore esercizio) ulteriore adattamento;
  \item senza allenamento: esercizio $\rightarrow$ debilitazione $\rightarrow$ ripristino, senza
    adattamento (``no allenamento'').
\end{itemize}

\subsection{Comando centrale nell'esercizio fisico}

Regolazione anticipatoria (\textit{feedforward}) a partenza dal \textbf{SNC} (corteccia cerebrale,
sistema limbico-ipotalamo), con due vie efferenti del sistema nervoso autonomo ortosimpatico
(SNA-OS):
\begin{itemize}
  \item \textbf{via colinergica}: SNC $\rightarrow$ SNA-OS colinergico $\rightarrow$ vasodilatazione
    (muscolo) $\rightarrow$ aumento del flusso;
  \item \textbf{via adrenergica}: SNC $\rightarrow$ SNA-OS adrenergico $\rightarrow$ cuore
    $\rightarrow$ aumento di gittata cardiaca e pressione ematica $\rightarrow$ (gittata) aumento
    del flusso.
\end{itemize}

\rubrica{Retroazioni (feedback)}
\begin{itemize}
  \item \textbf{movimento} $\rightarrow$ attivazione dei recettori muscolari ed articolari
    $\rightarrow$ produzione di metaboliti da lavoro $\rightarrow$ aumento del flusso; e segnale in
    ritorno verso il tronco encefalico e il SNC;
  \item \textbf{tronco encefalico} (centri di controllo cardiovascolari): riceve il feedback (dalla
    pressione ematica e dal movimento), agisce in senso inibitorio ($-$) sull'SNA-OS adrenergico e
    invia segnale eccitatorio ($+$) al SNC.
\end{itemize}

\subsection{Regolazione dei principali organi e sistemi durante l'esercizio}

\begin{table}[htbp]
  \centering
  \small
  \renewcommand{\arraystretch}{1.3}
  \begin{tabularx}{\textwidth}{|l|X|l|}
    \hline
    \textbf{Sistema} & \textbf{Risposta} & \textbf{Variabile regolata} \\
    \hline
    Circolatorio & Aumento della frequenza, aumento del flusso & Pressione ematica \\
    \hline
    Respiratorio & Aumento della frequenza & Ossigenazione \\
    \hline
    Rene & Riassorbimento di acqua e sali & Volume ematico \\
    \hline
    Fegato & Aumento dell'immissione di glucosio & Glicemia \\
    \hline
    Muscolo scheletrico & Contrazione-rilasciamento; movimento, aumento del carico & --- \\
    \hline
    Endocrino & Regolazione dei livelli ormonali & --- \\
    \hline
    Cute & Aumento della sudorazione & Temperatura corporea \\
    \hline
    Nervoso & Aumento/decremento dell'attività & --- \\
    \hline
    Digerente & Diminuzione dell'attività & --- \\
    \hline
  \end{tabularx}
  \caption{Risposta dei principali sistemi durante l'esercizio fisico e variabile regolata.}
  \label{tab:regolazione-esercizio}
\end{table}

% ---------------------------------------------------------------------------
\section{Integrazione tra sistema nervoso ed endocrino}

Le interazioni tra sistema nervoso ed endocrino mantengono l'omeostasi, soprattutto a fronte di
variazioni repentine dei parametri come nell'attività sportiva. Tutte le attività richiedono
un'integrazione coordinata dei sistemi biologici e fisiologici.
\begin{itemize}
  \item \textbf{sistema endocrino (S.E.)}: coordinamento fine delle risposte fisiologiche a un
    qualsiasi disturbo dell'equilibrio omeostatico;
  \item \textbf{sistema nervoso (S.N.)}: comunicazioni tra tessuti e sistemi dell'organismo.
\end{itemize}

\rubrica{Ruolo dell'ipotalamo}
\textbf{Ipotalamo}: produce il precursore dei peptidi oppioidi (endorfine) e dell'ACTH (ormone
adrenocorticotropo) $\rightarrow$ adattamenti comportamentali (paura, aggressività, curiosità,
vigilanza, dolore, ecc.).

\subsection{Bilancio idrico ed ADH nell'esercizio}

A seguito di un esercizio fisico intenso si ha sudorazione intensa per facilitare la dispersione di
calore; la perdita di acqua attiva gli osmocettori, che inducono la produzione dell'ormone
antidiuretico (ADH) da parte dell'ipotalamo.

\textbf{Sequenza} (attività fisica $\rightarrow$ sudorazione $\rightarrow$ perdita di $H_2O$):
\begin{enumerate}
  \item perdita della componente liquida del sangue $\rightarrow$ emoconcentrazione e aumento
    dell'osmolarità del sangue;
  \item l'aumento dell'osmolarità stimola gli \textbf{osmocettori ipotalamici} $\rightarrow$ si
    attiva il \textbf{centro della sete} (aumento dell'ingestione di acqua) e i neuroni ipotalamici
    producono \textbf{ADH}, che viene secreto dalla neuroipofisi;
  \item l'\textbf{ADH} agisce su dotti e tubuli renali $\rightarrow$ aumenta il riassorbimento di
    acqua $\rightarrow$ diminuisce l'osmolarità del sangue, ristabilendo l'osmolarità dei liquidi
    intra ed extracellulari.
\end{enumerate}

\subsection{Differenze tra controllo nervoso ed endocrino}

\begin{table}[htbp]
  \centering
  \small
  \renewcommand{\arraystretch}{1.3}
  \begin{tabularx}{\textwidth}{|l|X|X|}
    \hline
    & \textbf{Sistema nervoso} & \textbf{Sistema endocrino} \\
    \hline
    Trasmissione del segnale & Sinapsi & Sangue \\
    \hline
    Segnale trasmesso & Elettrochimico & Chimico \\
    \hline
    Sede della risposta & Localizzata & Diffusa \\
    \hline
    Velocità della risposta & Rapida & Lenta \\
    \hline
    Durata della risposta & Breve & Duratura \\
    \hline
    Processi controllati & Rapidi & Lenti \\
    \hline
  \end{tabularx}
  \caption{Differenze tra i controlli nervoso ed endocrino.}
  \label{tab:nervoso-endocrino}
\end{table}
